%% Generated by verification/make_figures.py from regular-graph-harmonic-quotient-examples.tex.
%% Figure fig:harmonic-quotient_d3reg10k+house (replaces figures/harmonic-quotient_d3reg10k+house-G2.tex).
%% Gamma_2 = d3reg10k: 10 vertices, 15 edges (E_2 in the example).
%% phi from the table at line 2311. Vertex x is filled with color hq<phi(x)>.
%% Layout: circle, cycle 0,1,...,9.
\begin{tikzpicture}[hqgraph]
\node[hqv, fill=hq2] (v0) at (0.000,1.770) {0};
\node[hqv, fill=hq2] (v1) at (-1.040,1.432) {1};
\node[hqv, fill=hq3] (v2) at (-1.683,0.547) {2};
\node[hqv, fill=hq0] (v3) at (-1.683,-0.547) {3};
\node[hqv, fill=hq4] (v4) at (-1.040,-1.432) {4};
\node[hqv, fill=hq1] (v5) at (-0.000,-1.770) {5};
\node[hqv, fill=hq1] (v6) at (1.040,-1.432) {6};
\node[hqv, fill=hq4] (v7) at (1.683,-0.547) {7};
\node[hqv, fill=hq0] (v8) at (1.683,0.547) {8};
\node[hqv, fill=hq3] (v9) at (1.040,1.432) {9};
\draw[hqe] (v0) -- (v1);
\draw[hqe] (v0) -- (v6);
\draw[hqe] (v0) -- (v9);
\draw[hqe] (v1) -- (v2);
\draw[hqe] (v1) -- (v5);
\draw[hqe] (v2) -- (v3);
\draw[hqe] (v2) -- (v4);
\draw[hqe] (v3) -- (v4);
\draw[hqe] (v3) -- (v8);
\draw[hqe] (v4) -- (v5);
\draw[hqe] (v5) -- (v6);
\draw[hqe] (v6) -- (v7);
\draw[hqe] (v7) -- (v8);
\draw[hqe] (v7) -- (v9);
\draw[hqe] (v8) -- (v9);
\end{tikzpicture}%
